基于CMS实验在四轻子末态测量希格斯玻色子产生截面

The measurement of Higgs production cross section in the four-lepton final state at CMS

CMS

Higgs, cross section, CMS

201609-202008

2020-08-25

One of the important properties of the Higgs boson is its model independent production cross section for a particular process, where the cross section illustrates the possibility of the occurrence of that process. 
Higgs cross section measurements are important to test the Standard Model predictions in the Higgs sector and can be used to constrain phase spaces of many theories beyond the Standard Model. Full Run2 data would allow sufficient statistics and has allowed to use finer bins than previous measurements on the same topic.

1) G. Aad et al., “Measurements of the Total and Differential Higgs Boson Production Cross Sections Combining the H →γγ and H →ZZ→4l Decay Channels at √s = 8 TeV with the ATLAS Detector”, Phys. Rev. Lett. 115(2014), p. 091801.

2) ATLAS collaboration, “Measurements of the total cross sections for Higgs boson production combining the H →γγ and  H→ZZ→4l decay channels at 7, 8 and 13 TeV center-of-mass energies with the ATLAS detector”, ATLAS Note ATLAS-CONF-2015-069 (2015), pp. 1–61.

3) Vardan Khachatryan et al. “Measurement of differential and integrated fiducial cross sections for Higgs boson production in the four-lepton decay channel in pp collisions at s = 7 and 8 TeV”. In: JHEP 04 (2016), p. 005.

4) Albert M Sirunyan et al. “Measurements of properties of the Higgs boson decaying into the four-lepton final state in pp collisions at √s =13 TeV”. In: JHEP 11 (2017), p. 047.

5) Measurements and interpretations of Higgs-boson fiducial cross sections in the diphoton decay channel using 139 fb ‘ ‘1 of pp collision data at √ s = 13 TeV with the ATLAS detector. Tech. rep. ATLAS-CONF-2019-029. Geneva: CERN, 2019.

6) Combined measurement of the total and differential cross sections in the H → γγ and the H → ZZ∗ → 4l decay channel7 at √s =13 TeV with the ATLAS detector. Tech. rep. ATLAS-CONF-2019-032. Geneva: CERN, 2019.

5) V. Khachatryan et al. “Measurement of differential cross sections for Higgs boson production in the diphoton decay channel in pp collisions at s = √8 TeV”. In: Eur. Phys. J. C 76 (2016), pp. 1–31.

The analysis on which the dissertation is built, was the first measurement with full Run 2 CMS data.
I performed Higgs inclusive and differential cross section measurements. Additionally, I performed muon object studies with full Run 2 data for this analysis.


本论文通过希格斯玻色子衰变到四轻子过程研究了希格斯粒子的性质。基于CMS探

测器收集的质心系能量为13 TeV的质子质子对撞数据,测量了希格斯至四轻子衰

变的基准截面,包括积分截面和微分截面。这些测量使用了CMS第二期运行获

取的全部数据,即质心系能量为13 TeV、积分亮度为137.1 fb −1 的数据,在测量

中希格斯质量取值为125.09 GeV。测量结果直接检验了标准模型希格斯相关的

在TeV能标上的QCD微扰计算。为了使所测量结果不依赖于希格斯的理论预测,

论文基于末态轻子运动学信息和事例拓扑结构定义了基准相空间,所有的截面

测量都限定与基准相空间。在质心系能量为13 TeV情况下,积分基准截面测量值

为2.73 (+0.30)(−0.29) fb。 另外针对多个对希格斯玻色子产生机制

敏感的运动学观测量,包括快度、四轻子体系的横动量、相关的多重喷注以及领头

喷注的横动量等进行了基准微分截面测量。测量结果对于检验标准模型希格斯粒子

的预测至关重要,同时可以用于约束许多超标准模型理论的相空间。在误差范围

内,积分截面和微分截面结果均与标准模型理论预言一致。




The measurement of properties of Higgs boson is presented in its four-lepton decay channel. The fiducial cross sections (integrated and differential) of the Higgs boson production via H → 4l decays (l = e, μ) are measured using the data collected with CMS detector in pp collisions at √s = 13 TeV . The measurements

use Higgs mass m H = 125.09 GeV and correspond to an integrated luminosity of 137.1 fb −1 at 13 TeV(Full Run 2 data). These measurements provide a direct test of perturbative QCD calculations in the Higgs sector of Standard Model at the TeV energy scale.

In order to make all the measurements independent of how Higgs is produced, all the cross sections are extracted within a fiducial phase space defined by the selections on lepton kinematics and event topology. The integrated fiducial cross sections are measured to be 2.73 (+0.30)

(−0.29) fb at 13 TeV. The differential fiducial cross sections are presented as a function of several kinematic

observables which are sensitive to the Higgs-boson production mechanisms: rapidity and transverse momentum of the four-lepton system, associated jet multiplicity and transverse momentum of the leading jet. These measurements are important to test the Standard Model predictions in the Higgs sector and can be

used to constrain phase spaces of many theories beyond the Standard Model.

All the integrated and differential measurements are compared with the theoretical predictions based on the Standard Model and found to be consistent within their uncertainties.


